%%%%%%%%%%%%%%%%%%%%%%%%%%%%%%%%%%%%%%%%%%%%%%%%%%%%%%%%%%
%%
%% MARK: - Chapter 0 — Preface
%%
%%%%%%%%%%%%%%%%%%%%%%%%%%%%%%%%%%%%%%%%%%%%%%%%%%%%%%%%%%

\chapter*{Preface}

For nearly a century,
quantum mechanics has rested upon two distinct conceptual pillars.
The first,
erected by Dirac and refined by Kostant and Souriau,
is the pillar of operators acting on a Hilbert space.
The second,
envisioned by Feynman,
is the pillar of the ``sum-over-histories,''
where the quantum amplitude is distilled from the space of all paths in configuration space.

While both approaches have achieved unparalleled empirical success,
the geometric bridge between them has remained elusive.
Standard geometric quantization has largely followed the Dirac program,
constructing principal bundles over the classical space of motions.
Yet these bundles often feel like external additions to the classical geometry,
their connection forms chosen by gauge rather than given by the physics.
On the other hand,
the Feynman path integral,
despite its profound physical intuition,
has remained a mathematical ghost,
trapped by the impossibility of defining a translation-invariant measure on the space of all paths.

This Memoir proposes a fundamental shift of perspective.
We do not sum over paths in configuration space;
we filter the space of \emph{finite fluctuations between complete classical motions}.
In the language of Souriau,
the classical system is described not by instantaneous snapshots
but by the symplectic space $(X,\omega)$ of entire dynamical trajectories.
A point in $X$ is not a state at an instant;
it is a complete motion.
Our ``paths'' are not trajectories of a particle;
they are \emph{deformations} connecting one classical dynamical history to another.
A path $\gamma \in \Paths(X)$ with endpoints $\gamma(0) = x$ and $\gamma(1) = x'$
represents a finite fluctuation from the motion $x$ to the motion $x'$.

The source of all quantum information is the space of these fluctuations.
By working within the framework of \emph{Diffeology},
we can treat the infinite dimensional space $\Paths(X)$ of all smooth fluctuations as a geometrically sound object.
Diffeology allows us to take quotients without losing smoothness,
and to handle spaces---such as irrational tori or leaf spaces of foliations---that lie beyond the reach of traditional manifolds.
We show that the quantum system is not a bundle to be added to the classical space,
but a groupoid to be distilled from the fluctuations of classical motions.

The core of our approach,
which we call \emph{Geometric Quantization by Paths} (GQbP),
consists of a rigorous filtration of this space of fluctuations.
We call $\Paths(X)$ the \emph{Quantum Broth}.
It is a primary geometric object,
equipped with a canonical $1$-form,
the \emph{Action Form} $\Komega$,
which serves as a universal,
intrinsic Lagrangian.
Even when the classical symplectic form $\omega$ is not exact on $X$---as it is for the sphere $S^2$, for example---it is always exact on the space of fluctuations.
The Action Form is its canonical primitive.

The process of quantization is then revealed to be a distillation.
The principal mathematical contribution of this Memoir is the construction of the prequantum groupoid $\TT_\omega$ as the quotient:
\[
  \TT_\omega = \Paths(X)/\!\sim_\omega,
  \]
where the equivalence relation $\sim_\omega$ identifies two fluctuations $\gamma$ and $\gamma'$
with the same endpoints whose concatenation encloses a symplectic area belonging to the period group $\Pomega$.%
\footnote{More precisely,
the equivalence relation is defined by the vanishing of the Chasles cocycle $\Phi$,
derived from the \emph{paths-primitive} $1$-form $\Komega$.}
The symplectic area enclosed by the concatenated fluctuation is its \emph{quantum phase}.
The filtration discards those fluctuations whose relative quantum phase vanishes modulo the periods,
isolating the constructive interference.
The result is the prequantum groupoid $\TT_\omega$,
whose space of morphisms $\cY$ carries the \emph{Prequantum $1$-form} $\blambda$,
which is left and right invariant and whose curvature encodes the classical form $\omega$.

A critical feature of this construction is its absolute canonicity.
From the classical form $\omega$ to the Action Form $\Komega$ on the Broth,
and finally to the prequantum $1$-form $\blambda$ on the groupoid,
every step is geometrically determined.
There is no arbitrary choice of primitive;
the prequantum structure is an inherent property of the parasymplectic geometry,%
\footnote{A \emph{parasymplectic} form is,
by definition,
any closed $2$-form.
No non-degeneracy condition is assumed.}
not a structure to be assembled through gauge choices.

We interpret this groupoid,
the pair $(\cY,\blambda)$,
as the \emph{Space of Quantum Motions}.
In this geometry,
the classical space $X$ survives as the \emph{Skeleton} of identity morphisms,
representing the ``null'' fluctuations.
Surrounding this skeleton is the \emph{Quantum Fog} of non-unit morphisms,
which constitute the channels of quantum transition.
These morphisms encode the physically meaningful fluctuations:
they are the geometric reification of Feynman's interfering alternatives,
transposed from configuration space to the space of motions.

This construction establishes a rigorous parallel between classical and quantum mechanics.
Just as the pair $(X, \omega)$ constitutes the classical space of motions,
the pair $(\cY, \blambda)$ constitutes the space of quantum motions.
The prequantum $1$-form $\blambda$ is the quantum counterpart to the symplectic form,
encoding the transition phases between classical histories.
This symmetry reveals that quantization is not a departure from the geometry of motions,
but its natural extension into the relational domain of the groupoid.

The construction yields two fundamental structural results:
\begin{enumerate}
  \item \emph{The Intrinsic Phase}:
  The quantum phase is not an external structure added to the classical space.
  It emerges intrinsically as the isotropy group of the groupoid,
  isomorphic to the \emph{Torus of Periods} $\Tomega = \RR/\Pomega$.
  This torus captures the symplectic area of loops and the geometry of the fundamental group.
  \item \emph{Faithful Symmetries}:
  The group of automorphisms of the quantum system,
  $\Aut(\cY,\blambda)$,
  is naturally isomorphic to the group of symmetries of the classical system,
  $\Diff(X, \omega)$.
  The symmetries lift fully and faithfully to the groupoid without the need for central extensions at the geometric level.
\end{enumerate}

The Memoir is structured to follow the logical maturation of this framework.

\textbf{Part I} establishes the construction for connected and simply connected spaces.
We show how the quantum phase emerges naturally as the isotropy group.
All fundamental constructions are exposed in detail and can be appreciated for their internal logic without the clutter of homotopic technicalities.

\textbf{Part II} extends the theory to arbitrary connected spaces.
We identify the obstructions arising from non-trivial homotopy and resolve them by defining the \emph{Total Group of Periods}.
This allows the GQbP framework to handle systems with complex homotopy,
such as the Aharonov-Bohm effect or irrational tori,
with the same geometric ease as the simply connected case.

\textbf{Part III} crosses the ``Threshold of Analysis.''
We transition from the geometry of the groupoid to the analysis of its observables by constructing the intrinsic convolution algebra of half-densities.
Through the \emph{Experimentum Crucis} of the harmonic oscillator,
we demonstrate that this framework is not a cosmetic re-description but an operative machine.
We derive the exact quantum spectrum,
including the zero-point energy $n\hbar/2$,
strictly from the transformation laws of the filtered Broth.
The Metaplectic Anomaly appears not as a correction to be made,
but as a geometric inevitability to be revealed.

\textbf{Part IV} applies the framework to quantum field theory.
We treat the Wess-Zumino-Witten model,
the $O(3)$ non-linear sigma model on loop space,
and the Virasoro model.
These examples demonstrate that the GQbP framework scales from finite-dimensional systems to the infinite-dimensional geometry of loop spaces and conformal field theory,
providing a unified geometric language for quantum phenomena across scales.

This framework demonstrates that geometric quantization,
when properly understood,
is not incompatible with Feynman's path integral;
it is its direct geometric descendant---transposed to the space of motions and refined by the language of diffeology.
Diffeology does not force a choice between the two pillars of quantum mechanics.
Instead,
it preserves both within a single \emph{triptych} of spaces:
the \emph{Quantum Broth} of fluctuations,
the \emph{Prequantum Groupoid} of distilled quantum motions,
and the classical \emph{Skeleton} of complete histories.
Each level of this hierarchy is a rigorously structured diffeological space,
ensuring that the transition from fluctuations to operators remains geometrically transparent.

Why this filtration should recover the correct quantum system is a question geometry alone cannot answer.
What geometry provides is the canonical mechanism:
given the classical space of motions and its symplectic form,
the prequantum groupoid is the unique object that emerges from identifying fluctuations whose relative quantum phase belongs to the period group.
That this groupoid coincides with the Kostant-Souriau framework for manifolds,
and extends it faithfully to singular and infinite-dimensional spaces,
is the mathematical content of this Memoir.

Ultimately,
this work shows that geometric quantization is the logical conclusion of the Feynman vision---when that vision is lifted from configuration space to the space of motions and expressed in the precise language of diffeology.
The prequantum groupoid stands as the bridge,
poised exactly at the point where the geometry of fluctuations yields the analysis of quantum states.

\vspace{2em}
\begin{flushright}
  \emph{``Quoi de neuf en diffeolog'y ?''} \\
  --- Jean-Marie Souriau
\end{flushright}
